% !TeX program = xelatex
% ECSE 429 Project Part A -- Group 3
% Working scaffold; grey prompts need evidence before submission.
% Adapted from ECSE 211 Profile 1. Target: 5--10 pages.
% Manual page breaks group the draft sections. Reflow as content is added.
\documentclass[12pt,letterpaper]{article}
% ECSE 429 adaptation of the ECSE 211 Profile 1 visual system.
\usepackage{iftex}
\ifPDFTeX
  \usepackage[T1]{fontenc}
  \usepackage{tgtermes}
\else
  \usepackage[no-math]{fontspec}
  \setmainfont{Times New Roman.ttf}[
    Path=fonts/,
    BoldFont=Times New Roman Bold.ttf,
    ItalicFont=Times New Roman Italic.ttf,
    BoldItalicFont=Times New Roman Bold Italic.ttf]
\fi
\usepackage[letterpaper,margin=.8in]{geometry}
\usepackage{xcolor,graphicx,amsmath,newtxmath}
\SetSymbolFont{operators}{normal}{OT1}{ntxtlf}{m}{n}
\SetSymbolFont{operators}{bold}{OT1}{ntxtlf}{b}{n}
\SetMathAlphabet{\mathrm}{normal}{OT1}{ntxtlf}{m}{n}
\usepackage{tikz}
\usetikzlibrary{arrows.meta,shapes.geometric,positioning,calc}
\usepackage{array,booktabs,tabularx,enumitem}
\usepackage{caption,fancyhdr,listings}
\usepackage{xurl}
\definecolor{Heading}{RGB}{31,42,60}
\definecolor{Rule}{RGB}{154,167,189}
\definecolor{Caption}{RGB}{90,107,133}
\definecolor{DiagramBlue}{RGB}{49,77,121}
\definecolor{DiagramFill}{RGB}{240,243,247}
\definecolor{DraftInk}{RGB}{92,102,116}
\usepackage[unicode,colorlinks=true,linkcolor=Heading,urlcolor=Heading,
  citecolor=Heading,filecolor=Heading,runcolor=Heading,
  pdfdisplaydoctitle=true]{hyperref}
\hypersetup{
  pdftitle={ECSE 429 -- Project Part A Report},
  pdfauthor={James Helou; Jeffrey Lim; Mahan Maalekian; Matthew Petruzziello},
  pdfsubject={Group 3 -- exploratory testing of the Todo Manager REST API}}

\newcommand{\BodyFont}{\fontsize{12bp}{13.8bp}\selectfont}
\newcommand{\TableFont}{\fontsize{10.5bp}{12.2bp}\selectfont}
% Keep headings with their next paragraph without an extra package dependency.
\newcommand{\Needspace}[1]{%
  \par\begingroup\dimen0=#1\relax
  \dimen2=\pagegoal\advance\dimen2 by -\pagetotal
  \ifdim\dimen2<\dimen0\newpage\fi\endgroup}
\setlength{\parindent}{0pt}
\setlength{\parskip}{4bp}
\setlength{\topskip}{0pt}
\setlength{\emergencystretch}{1em}
\raggedright
\AtBeginDocument{\BodyFont}
\pagestyle{fancy}
\fancyhf{}
\renewcommand{\headrulewidth}{0pt}
\renewcommand{\footrulewidth}{0pt}
\fancyfoot[L]{\fontsize{9bp}{11bp}\selectfont\color{Caption}ECSE 429 \quad Group 3}
\fancyfoot[R]{\fontsize{9bp}{11bp}\selectfont\color{Caption}\thepage}

% Retain the ECSE 211 cover's centred blocks and measured baselines.
\newcommand{\CoverLine}[3]{%
  \node[anchor=base,inner sep=0pt,font=\fontsize{#2bp}{#2bp}\selectfont]
    at ([yshift=-#1bp]current page.north) {#3};}
\newcommand{\ReportCover}{%
  \thispagestyle{empty}
  \begin{tikzpicture}[remember picture,overlay]
    \CoverLine{100.14}{16}{\textbf{McGill University}}
    \CoverLine{114.8}{12}{Department of Electrical and Computer Engineering}
    \CoverLine{128.6}{12}{\CourseTitle}
    \CoverLine{212.53}{28}{\textbf{\ProjectTitle}}
    \CoverLine{233.52}{16}{\ProjectSubtitle}
    \CoverLine{254.5}{12}{Report}
    \draw[Rule,line width=.3bp]
      ([xshift=57.6bp,yshift=-271.5bp]current page.north west)
      -- ([xshift=-57.6bp,yshift=-271.5bp]current page.north east);
    \CoverLine{310.65}{14}{\textbf{Group \GroupNumber}}
    \node[anchor=north,inner sep=0pt] at
      ([yshift=-329.68bp]current page.north) {%
      \begin{tabular}{@{}l@{\hspace{1.8em}}r@{}}
        \TeamRoster
      \end{tabular}};
    \CoverLine{472.47}{12}{Submitted \SubmissionDate}
  \end{tikzpicture}
  \null\newpage}

\newcommand{\ReportSection}[1]{%
  \par\Needspace{55bp}\addvspace{6bp}
  {\color{Heading}\fontsize{16bp}{18.4bp}\selectfont\bfseries
   \rule{0pt}{14.94bp}%
   \rlap{\raisebox{-8.26bp}[0pt][0pt]{\color{Rule}\rule{\linewidth}{.3bp}}}%
   #1\par}
  \nobreak\vspace{-2bp}\hrule width 0pt height .3bp
  \nobreak\vspace{5bp}}
\newcommand{\ReportSubsection}[1]{%
  \par\Needspace{40bp}\addvspace{5bp}
  {\color{Heading}\fontsize{13bp}{15bp}\selectfont\bfseries
   \rule{0pt}{12.14bp}#1\par}
  \nobreak\vspace{1bp}}
\newcommand{\RunIn}[1]{\textbf{#1:}\ }
\newcommand{\NewContentPage}{\newpage\BodyFont}
\newcommand{\Draft}[1]{\textcolor{DraftInk}{\textit{[#1]}}}
\newcommand{\Prompt}[1]{\par{\color{DraftInk}\textit{[#1]}}\par}
\newcommand{\DraftNotice}{%
  {\fontsize{10bp}{12bp}\selectfont\color{Caption}%
  WORKING OUTLINE \quad Grey bracketed text and figure panels are drafting
  placeholders, not reported findings.\par}}
\newcommand{\Path}[1]{\texttt{\fontsize{10.5bp}{12.2bp}\selectfont\detokenize{#1}}}
\newcolumntype{Y}{>{\raggedright\arraybackslash}X}
\newcolumntype{P}[1]{>{\raggedright\arraybackslash}p{#1}}
\newenvironment{ReportTable}[1]{%
  \par\begingroup\TableFont
  \setlength{\tabcolsep}{5bp}\renewcommand{\arraystretch}{1.16}
  \tabularx{\linewidth}{#1}%
}{\endtabularx\par\endgroup\vspace{2bp}}
\newenvironment{Charter}{%
  \par\begingroup\setlength{\leftskip}{12bp}\setlength{\rightskip}{12bp}
  \color{Heading}\itshape
}{\par\endgroup}

\DeclareCaptionFont{captionblue}{\color{Caption}}
\captionsetup[figure]{hypcap=false,font={small,captionblue},labelfont=bf,textfont=normalfont,
  labelsep=period,justification=centering,skip=5bp}
\newcommand{\FigurePlaceholder}[5]{%
  \par\begingroup\centering
  \begin{tikzpicture}
    \node[draw=Rule,fill=DiagramFill,line width=.55bp,rounded corners=4bp,
      minimum height=#1,text width=\dimexpr\linewidth-26bp\relax,
      align=center,inner sep=10bp] {%
      {\color{DiagramBlue}\fontsize{11bp}{13bp}\selectfont\bfseries #2}\\[5bp]
      {\color{Caption}\fontsize{10bp}{12bp}\selectfont DIAGRAM PLACEHOLDER}\\[4bp]
      {\color{DraftInk}\fontsize{10bp}{12bp}\selectfont #3}};
  \end{tikzpicture}
  \captionof{figure}{#4}\label{#5}
  \par\endgroup\vspace{3bp}}
% Styles available for the later diagram-authoring pass.
\tikzset{
  process/.style={draw=DiagramBlue,fill=DiagramFill,line width=.5bp,
    rounded corners=5bp,align=center,inner sep=5bp,
    font=\sffamily\fontsize{9bp}{10.5bp}\selectfont},
  terminal/.style={process,fill=DiagramBlue,text=white},
  decision/.style={diamond,aspect=2.8,draw=DiagramBlue,fill=DiagramFill,
    align=center,inner sep=2bp,font=\sffamily\fontsize{9bp}{10.5bp}\selectfont},
  arrow/.style={-{Stealth[length=4bp,width=3bp]},DiagramBlue,line width=.5bp},
  branch/.style={font=\sffamily\fontsize{8bp}{9bp}\selectfont,text=black}
}
\lstset{
  basicstyle=\ttfamily\fontsize{9.5bp}{11.5bp}\selectfont,
  columns=fullflexible,keepspaces=true,breaklines=true,
  frame=single,framerule=.35bp,rulecolor=\color{Rule},
  backgroundcolor=\color{DiagramFill},framesep=6bp,
  aboveskip=5bp,belowskip=6bp,showstringspaces=false,
  xleftmargin=7bp,xrightmargin=7bp}


% -------------------------- Cover metadata --------------------------
\newcommand{\CourseTitle}{ECSE 429: Software Validation}
\newcommand{\ProjectTitle}{Project Part A}
\newcommand{\ProjectSubtitle}{Exploratory Testing of REST API}
\newcommand{\SubmissionDate}{October 20, 2026}
\newcommand{\GroupNumber}{3}
% All IDs verified against page 1 of the current session-note PDFs.
\newcommand{\TeamRoster}{%
  James Helou & 261177052\\
  Jeffrey Lim & 261146821\\
  Mahan Maalekian & 261175419\\
  Matthew Petruzziello & 261140144}

% Shared artifact links keep the submission table and references consistent.
% Paths are relative to this report in the ZIP's Report & Documents/ folder.
\newcommand{\MainReportLink}{%
  \href{ECSE429_PartA_Group3_Report.pdf}%
  {\Path{ECSE429_PartA_Group3_Report.pdf}}}
\newcommand{\BugSummaryLink}{%
  \href{Bug Summary - ECSE 429 Part A - Group 3.pdf}%
  {\Path{Bug Summary - ECSE 429 Part A - Group 3.pdf}}}
\newcommand{\CapabilityInventoryLink}{%
  \href{ECSE429_PartA_Group3_CapabilityInventory.pdf}%
  {\Path{ECSE429_PartA_Group3_CapabilityInventory.pdf}}}
\newcommand{\UnitTestVideoLink}{%
  \href{../Videos/ECSE429_PartA_Group3_UnitTests.mp4}%
  {\Path{Videos/ECSE429_PartA_Group3_UnitTests.mp4}}}
\newcommand{\NotesRef}[1]{%
  \hyperref[ref:session-#1]{S#1 session notes~[\ref*{ref:session-#1}]}}
\newcommand{\BugSummaryRef}{%
  \hyperref[ref:bug-summary]{bug summary~[\ref*{ref:bug-summary}]}}
\newcommand{\SessionNotesFolder}{%
  \href{../Session\%20Notes/}{\Path{Session Notes/}}}
\newcommand{\SessionScriptsFolder}{%
  \href{../Session\%20Scripts/}{\Path{Session Scripts/}}}
\newcommand{\SuiteFolder}{\href{../testsuite/}{\Path{testsuite/}}}
\newcommand{\ArtifactFile}[2]{%
  \href{#1}{{\fontsize{10.5bp}{12.2bp}\selectfont\nolinkurl{#2}}}}

\begin{document}
\ReportCover
\setcounter{page}{1}

% ===================== BODY PAGE 1: DELIVERABLES =====================
\ReportSection{1. Summary of Deliverables}

\ReportSubsection{1.1 Project Objective and Scope}
We explored Todo Manager 1.5.5, a REST API for managing todos, projects, categories,
and the relationships between them. We used eight exploratory sessions, each
timeboxed at 45 minutes, and built Java unit tests from the capabilities and
concerns we identified. Our exploration found inconsistencies in relationship
queries, deletion cleanup, and malformed-input handling. The session notes and
bug summary record these findings, along with potential bugs that need retesting
or clarification of the expected behaviour before we can classify them as defects.

\ReportSubsection{1.2 Submission Package}
\begin{ReportTable}{@{}P{.25\linewidth}Y@{}}
  \toprule
  \textbf{Deliverable} & \textbf{Location and role}\\
  \midrule
  Main report PDF &
  \MainReportLink\ in \Path{Report & Documents/}.
  Summarises our exploration, test suite, submission package, results, and AI use.\\
  Session-note PDFs &
  \SessionNotesFolder:
  eight PDFs for S1--S8, containing observations, findings, concerns, and new test ideas.\\
  Scripts and evidence &
  \SessionScriptsFolder:
  seven Python demonstrations for S1, S2, and S4--S8.
  S3's Postman requests are recorded in its session notes.\\
  Unit-test suite &
  \SuiteFolder:
  our Java tests using JUnit 5 and REST Assured, with the Maven configuration
  and wrapper needed to run them.\\
  Bug summary &
  \BugSummaryLink\ in \Path{Report & Documents/}.
  Entries BR01--BR33 describe findings, impact, and reproduction steps.
  BR30--BR33 flag potential bugs needing retesting or clearer expected behaviour.\\
  Capability inventory &
  \CapabilityInventoryLink\ in \Path{Report & Documents/}.
  Maps documented and discovered capabilities to session evidence,
  demonstrations, and Java test modules.\\
  Unit-test video &
  \UnitTestVideoLink.
  Demonstrates suite execution and pseudorandom test ordering.\\
  \bottomrule
\end{ReportTable}

\ReportSubsection{1.3 Reading the Evidence}
We use session and observation numbers to connect our findings to the requests
or scripts, bug-summary entries, and corresponding unit tests. This report
summarises that work, while the separate notes, scripts, and bug records provide
the detailed steps and observations.

\ReportSubsection{1.4 Team Participation}
James handled the project setup and S4, S7, and S8. Jeffrey completed S1 and S2
and wrote the todo tests. Mahan completed S3 and wrote the project tests.
Matthew completed S5 and S6, wrote the category tests, and was responsible for
the unit-test video. We all worked on the report together.

% ===================== BODY PAGE 2: EXPLORATION METHOD ================
\NewContentPage
\ReportSection{2. Exploratory Testing Method}

\ReportSubsection{2.1 Exploratory Testing Charter}
\begin{Charter}
Identify capabilities and areas of potential instability of the
\textquotedblleft rest api todo list manager\textquotedblright.

Identify documented and undocumented
\textquotedblleft rest api todo list manager\textquotedblright\ capabilities.

For each capability, create a script or small program to demonstrate
the capability.

Exercise each capability identified with data typical to the intended
use of the application.
\end{Charter}

\RunIn{Session approach}
We divided the exploration into sessions on resources and their relationships.
Each session PDF records one participant and a 45-minute timebox. Across the
sessions, we tried typical inputs as well as omitted fields, invalid IDs,
different request formats, and different operation orders.

We compared responses with the API documentation~[\ref{ref:api}] and used GET
requests to check what was stored. In S1, we checked that full updates worked
before comparing POST and PUT requests that left fields out
(\NotesRef{1}, observations 4--7, pp.~2--3). For relationships, we checked the links
from both objects. In S6, we compared plain XML, JSON, and XML CDATA for the same
target, removing the link after each successful request so that the next format
was tested on an unlinked pair (\NotesRef{6}, observation 14, p.~4).

Our notes record what we tried and expected, what we observed, the capabilities
we exercised, the risks we identified, and the decisions we made. We then
reviewed our findings and follow-up ideas.
Figure~\ref{fig:session-evidence} reads from top to bottom: the charter guides the
sessions, notes and demonstrations record the evidence, and selected concerns
and findings inform bug records and tests.

\begin{center}
  \includegraphics[width=\linewidth]{images/session-evidence-traceability.png}
  \captionof{figure}{How exploration informs supporting records and tests.
  Arrow labels describe each relationship. BR30--BR33 are potential bugs
  needing retesting or clearer expected behaviour.}
  \label{fig:session-evidence}
\end{center}

\NewContentPage
\ReportSubsection{2.2 Session Coverage and Participation}
% Keep the final sentence on this page without changing the text or font size.
\enlargethispage{\baselineskip}
We used eight sessions to explore todos, projects, categories, and their
interoperability. S1--S6 covered creating, reading, updating, and deleting
resources (CRUD), along with the routes used to link them. S7 examined deletion
across linked objects, while S8 explored malformed todo requests and invalid
operations.

% Links retain the exact filenames in the current submission, including S3's spelling.
\begin{ReportTable}{@{}P{.05\linewidth}P{.22\linewidth}YP{.13\linewidth}P{.14\linewidth}@{}}
  \toprule
  \textbf{ID} & \textbf{Participant} & \textbf{Focus} & \textbf{Session date} & \textbf{Note reference}\\
  \midrule
  S1 & Jeffrey Lim & Todo CRUD, field validation, JSON/XML, and invalid input.
     & 17 Oct 2026 & \NotesRef{1}, pp.~1--5\\
  S2 & Jeffrey Lim & Todo links to projects and categories; reads, writes, and unlinking.
     & 17 Oct 2026 & \NotesRef{2}, pp.~1--6\\
  S3 & Mahan Maalekian & Project CRUD, validation, task linking, and deletion.
     & 18 Oct 2026 & \NotesRef{3}, pp.~1--3\\
  S4 & James Helou & Project task and category links; queries, linking, and unlinking.
     & 12 Oct 2026 & \NotesRef{4}, pp.~1--4\\
  S5 & Matthew Petruzziello & Category CRUD, validation, filters, formats, and routing.
     & 12 Oct 2026 & \NotesRef{5}, pp.~1--8\\
  S6 & Matthew Petruzziello & Category links to todos and projects; formats and deletion cleanup.
     & 18 Oct 2026 & \NotesRef{6}, pp.~1--7\\
  S7 & James Helou & Links across resources and deletion of related objects.
     & 12 Oct 2026 & \NotesRef{7}, pp.~1--3\\
  S8 & James Helou & Malformed JSON/XML todo requests, unsupported verbs, and invalid operations.
     & 12 Oct 2026 & \NotesRef{8}, pp.~1--4\\
  \bottomrule
\end{ReportTable}

\RunIn{Reading the session records}
The first page of each PDF lists the participant's full name, student ID, and
student email address, along with the charter, scope, solo participation, and
timebox. The table lists the dates when the sessions took place.
Each note reference points to a separate PDF in
\SessionNotesFolder\ in the submission ZIP. Section~8 gives the exact filenames;
page references count from the first page of each PDF.

\RunIn{Demonstrations and follow-up}
We supplied seven Python demonstrations in \SessionScriptsFolder\ for S1, S2,
and S4--S8. For S3, the Postman requests are recorded in \NotesRef{3}
(pp.~1--3). NOTE: no separate collection or request export is included.

We collected our findings and concerns in the \BugSummaryRef,
supplied as \BugSummaryLink\ in \Path{Report & Documents/}.
It contains 33 records. BR30--BR33 (p.~6) are potential bugs that need further
checks before we can classify them as defects. BR30 needs clarification of what
the documentation allows. BR31--BR32 need reads before and after PUT to check
which fields or links remained stored, and BR33 needs retesting to explain
why the sessions recorded different results.

We used selected observations to develop tests in \SuiteFolder.
For example, the XML comparison in \NotesRef{6}, observation 14 (p.~4),
appears as BR06 in the \BugSummaryRef\ (p.~2).
The category--project
\href{../testsuite/src/test/java/ca/mcgill/ecse429/group3/categories/item/projects/PostXmlExpectedTest.java}{XML expected-behaviour}
and
\href{../testsuite/src/test/java/ca/mcgill/ecse429/group3/categories/item/projects/PostXmlObservedTest.java}{observed-behaviour tests}
check the documented and observed results separately.

% ===================== BODY PAGE 3: EXPLORATORY FINDINGS ===============
\NewContentPage
% Findings are supported by the session PDFs and scripts.
\ReportSection{3. Findings of Exploratory Testing}
% Check the requested page budget during compilation without printing extra text.
\edef\FindingsFirstPage{\number\value{page}}

\ReportSubsection{3.1 What We Learned in Each Session}
Ordinary CRUD generally worked. However, relationship reads, deletion, and parsing
needed closer checks. Numbers identify observations in the linked notes.
Follow-ups reflect session ideas and remaining checks. Some were explored later.
\begin{ReportTable}{@{}P{.06\linewidth}YYY@{}}
  \toprule
  \textbf{ID} & \textbf{Findings / capabilities} & \textbf{Concerns / risks} & \textbf{New testing ideas}\\
  \midrule
  \hyperref[ref:session-1]{S1 [\ref*{ref:session-1}]} &
  JSON/XML reads worked. POST kept omitted fields. PUT reset them (7, 9). &
  Boolean output strings cannot be reused directly as input (2, 8). &
  Compare null/omitted XML fields and allocation after other failures or restart.\\
  \hyperref[ref:session-2]{S2 [\ref*{ref:session-2}]} &
  Project links worked both ways. Selective unlinking preserved objects (2, 8--9). &
  Missing-todo reads returned unrelated targets. Category links were one-way (5, 7). &
  Try negative/deleted todo IDs, XML link bodies, and reverse-side category unlinking.\\
  \hyperref[ref:session-3]{S3 [\ref*{ref:session-3}]} &
  Project GET/DELETE returned 200. POST \texttt{\{\}} created a blank project (1--2, 7). &
  PUT storage is unconfirmed. Task-link rejection conflicts with S4 (4--5). &
  Replay PUT with before/after GETs. Vary XML, filters, and text length.\\
  \hyperref[ref:session-4]{S4 [\ref*{ref:session-4}]} &
  Task linking was reciprocal. Unlinking kept the todo (3, 8). &
  Missing-project reads returned both projects' tasks (6). &
  Try negative/large project IDs, ID-plus-title precedence, and links after restart.\\
  \hyperref[ref:session-5]{S5 [\ref*{ref:session-5}]} &
  JSON/XML creation and exact, case-sensitive filters worked (2, 10). &
  Encoded-slash ID returned the collection. Weighted Accept returned 404 (9, 20). &
  Vary path encodings, Unicode filters, array coercion, and item/HEAD Accept.\\
  \hyperref[ref:session-6]{S6 [\ref*{ref:session-6}]} &
  String-ID linking and selective unlinking preserved shared targets (2--8). &
  Missing-parent reads included unrelated targets. Plain XML IDs failed unlike JSON/CDATA (9, 14). &
  Vary XML ID whitespace/escapes and deletion/unlink order with multiple links.\\
  \hyperref[ref:session-7]{S7 [\ref*{ref:session-7}]} &
  Todo/project deletion cleaned task links and preserved other objects (1--2). &
  A deleted category remained in a todo's list. Deleted-project task reads used the wrong key (3--4). &
  Create a replacement category. Delete a todo linked to both project and category.\\
  \hyperref[ref:session-8]{S8 [\ref*{ref:session-8}]} &
  Valid JSON/XML created todos. Malformed-input results varied by syntax (1--7). &
  Unquoted keys succeeded. JSON parser errors exposed classes and mismatched media types (2--5). &
  Try long titles, deep nesting, missing Content-Type, and XML Accept on errors.\\
  \bottomrule
\end{ReportTable}

\ReportSubsection{3.2 Capabilities and Demonstration Scripts}
Our
\hyperref[ref:capabilities]{capability inventory~[\ref*{ref:capabilities}]}
maps the 127 method/route contracts from 19 Swagger route templates to Java modules
for Section~4. It also records discovered behaviours, script inputs, and effects.
The API reference~[\ref{ref:api}] provides \Path{/docs} and \Path{/docs/swagger}
for documentation discovery. Examples follow.

\begin{ReportTable}{@{}P{.25\linewidth}P{.18\linewidth}Y@{}}
  \toprule
  \textbf{Capability / API} & \textbf{Status and source} & \textbf{Typical-data demonstration}\\
  \midrule
  POST \Path{/todos} &
  Documented in the API reference. &
  \Path{S1_todo_crud.py}, C2--C2b (S1, 2):
  Create with ``Buy groceries'', ``Milk and bread'', and false.
  Then GET the returned ID to check stored values.\\
  POST \Path{/categories/:id/todos} &
  Documented in the API reference. &
  \Path{S6_category_relationships.py}, block 2 (S6, 2):
  Send the returned todo ID as a JSON string. The category lists the unchanged
  target without creating another todo.\\
  POST \Path{/projects/:id/tasks}, title-only body &
  Undocumented creation behaviour (S4, 7). &
  \Path{S4_project_relationships.py}, E8--E8b:
  Omit the ID and supply a title. A new todo is created and linked.
  Compare the todo collection before and after.\\
  \bottomrule
\end{ReportTable}
The inventory distinguishes documented capabilities, additional behaviours, and
Java follow-up checks. However, a source mapping does not establish a passing test.

\ReportSubsection{3.3 Tester Decisions and Areas of Risk}
\RunIn{Use populated controls}
An empty category list did not show whether the project existed
(\NotesRef{4}, 2, p.~1). Populated parents exposed unrelated targets. Therefore,
we kept different linked lists and direct GET controls
(\NotesRef{6}, 9, p.~3, and \Path{S6_category_relationships.py}, block 9).
Status alone would have missed the wrong data.

\RunIn{Check deletion from both sides}
We restarted S7's scenarios to isolate deletions
(\NotesRef{7}, 6, p.~2, and \Path{S7_interoperability_delete.py}, \Path{restart()}).
In S6, unlinking one direction left the reverse link intact. Therefore, we decided
to unlink both directions before deleting temporary objects
(\NotesRef{6}, 7, p.~3, and \Path{S6_category_relationships.py}, block 7).

\RunIn{Change one input at a time}
We compared plain XML, JSON, and CDATA for the same ID. We unlinked between successes
(\NotesRef{6}, 14, p.~4, and
\Path{S6_category_relationships.py}, block 12).
Only plain XML failed. Therefore, we kept separate format checks. The cause remains unknown.

\ReportSubsection{3.4 Confirmed Bugs and Open Questions}
\RunIn{Confirmed discrepancies}
The \BugSummaryRef\ gives reproduction steps. Its most consequential discrepancies
are below. Page numbers refer to that PDF.
\begin{ReportTable}{@{}P{.16\linewidth}P{.34\linewidth}Y@{}}
  \toprule
  \textbf{Record / page} & \textbf{Expected behaviour} & \textbf{What we observed and why it matters}\\
  \midrule
  BR01, p.~1 & Relationship reads stay scoped to their named parent. &
  Missing parents returned 200 and unrelated targets. The exact missing-parent
  status is unspecified. Returning another parent's data is the discrepancy.\\
  BR02, p.~1 & A project's task response uses the \Path{todos} field. &
  After deletion, \Path{projects} replaced \Path{todos}. Therefore, clients could miss the list.\\
  BR04--BR05, p.~1 & Deleted objects stop appearing as live related records. &
  Direct GET returned 404. However, one-way links still returned full deleted records.
  BR05 also showed the deleted target beside its replacement.\\
  BR06, p.~2 & Equivalent JSON/XML IDs identify the same existing target. &
  Plain XML returned 404. However, JSON strings and CDATA returned 201 for the same ID.\\
  BR12, p.~3 & Swagger specifies 405 for GET/HEAD/POST on individual link paths. &
  Existing category links returned 404, making an unsupported method look like
  a missing link.\\
  BR14--BR15, p.~3 & A JSON-acceptable read works. An item lookup stays scoped. &
  An Accept list allowing JSON returned 404. JSON alone returned 200.
  \texttt{/categories/\%2F} returned the collection.\\
  \bottomrule
\end{ReportTable}
Some JSON parse errors exposed Java classes and returned JSON-shaped bodies labelled
\Path{text/html} (BR07--BR08, p.~2).

\RunIn{Other observations and documentation gaps}
The 33 records are not all confirmed defects. The documentation does not clearly
require category links to work both ways (BR03, p.~1) or specify what repeated
link POSTs should return (BR16, p.~3). Inline creation adds a capability
(BR17, p.~3). Optional titles allow blank projects (BR18, p.~4). Numeric IDs,
filters, omitted PUT fields, and null amendments need clearer rules
(BR19--BR23, p.~4). BR29 (p.~5) is a documentation gap: category relationship
POSTs safely return 404 for a missing parent, but Swagger omits that response.

\RunIn{Open questions and untested concerns}
BR30--BR33 (p.~6) need confirmation. Supplied-ID rejection may be valid. Project
PUT needs before/after scalar and relationship reads. Additionally, S3's task-link 400 conflicts
with S4's 201. Encoded-slash writes remain untested (\NotesRef{5}, 9, p.~4).
\par

% ===================== BODY PAGE 4: SUITE STRUCTURE ===================
\NewContentPage
% Check after flushing Section 3, so a pending page break cannot hide overflow.
\ifnum\value{page}>\numexpr\FindingsFirstPage+2\relax
  \PackageError{report}{Section 3 exceeds the requested two pages}%
    {Shorten Section 3 before recompiling.}
\fi
\ReportSection{4. Structure of the Unit-Test Suite}
\edef\SuiteFirstPage{\number\value{page}}
\newcommand{\TestSource}[1]{%
  \href{../testsuite/src/test/java/ca/mcgill/ecse429/group3/#1}{\Path{#1}}}

\begin{center}
  \includegraphics[width=\linewidth]{images/unit-test-architecture.png}
  \captionof{figure}{Suite dependencies and report output. Domain and shutdown
  tests use separate, externally started Todo Manager processes. The JDK client is used for
  category POST controls that omit Content-Type.}
  \label{fig:test-architecture}
\end{center}

\ReportSubsection{4.1 Module Organisation and API Coverage}
We target Java 17 and use the open-source JUnit Jupiter framework (5.11.3) for one suite.
REST Assured 5.5.0 sends HTTP requests to the running application, and Maven
Surefire 3.5.2 runs the tests and writes the results (Figure~\ref{fig:test-architecture}).
All Java paths below start at \Path{testsuite/src/test/java/ca/mcgill/ecse429/group3/}.

\RunIn{Modules}
Our packages follow the routes. For example, \Path{todos/allTodos/} covers
\Path{/todos}, \Path{projects/item/tasks/} covers a project's task list, and
\Path{categories/item/projects/item/} covers an individual category--project link.
Each verb has its own module, such as \Path{GetTest.java} or \Path{DeleteTest.java}.
Our \hyperref[ref:capabilities]{capability inventory~[\ref*{ref:capabilities}]}
lists all 127 Swagger method--route contracts, their test packages and verb
modules, and the additional behaviours found during exploration.
These include inline creation, encoded-slash routing, TRACE, and BREW.
Shutdown has separate \Path{shutdown/GetExpectedTest.java} and
\Path{shutdown/GetObservedTest.java} modules. We select the
\Path{isolated-shutdown} profile and one fully qualified \Path{-Dtest} class
at a time.

\RunIn{Expected and observed behaviour}
We keep discrepancies in separate modules so that a working alternative does
not hide a failed expectation. The XML comparison in \NotesRef{6}
(observation 14, p.~4) is BR06 in the \BugSummaryRef\ (p.~2).
The category--project
\TestSource{categories/item/projects/PostXmlExpectedTest.java} expects plain XML
to link the same target as JSON and CDATA, returning 201. The separate
\TestSource{categories/item/projects/PostXmlObservedTest.java} checks the
observed 404 with unchanged state, followed by successful JSON and CDATA links.
We unlink between formats and check that the original state is restored, so an
existing link cannot explain the next result. A failed expectation remains an
assertion failure in Surefire.

\NewContentPage
\ReportSubsection{4.2 Test Lifecycle and Correctness}
\RunIn{Readiness and initial conditions}
\TestSource{utilities/ProjectTestSetup.java} reads all three collections in
\Path{@BeforeEach} and saves their records and relationship references through
\Path{ApiStateSnapshot}. A missing service fails setup.
\Path{SetupTestUtilities} prepares and tracks temporary records and links, so
tests do not depend on data left by an earlier test.

\RunIn{Assertions and side effects}
We check status codes against the API documentation where specified, then check
stored fields, collection membership, and links from the relevant directions.
\Path{ApiStateSnapshot.assertUnchangedExcept} excludes the records whose changes
are checked separately and compares the remaining visible data. The comparison
ignores list order while retaining duplicates and scalar types. Therefore, an
extra relationship or a changed unrelated record still causes a failure.
The BR06 modules above show this check alongside their response assertions.

\RunIn{Restoration after failure}
\Path{@AfterEach} runs cleanup and the initial-state comparison inside
\Path{assertAll}, including after an assertion fails. Owned links are removed
before owned records, in reverse order. Cleanup accepts 200 or 404, so an object
already deleted by the test does not prevent the remaining cleanup attempts.
The final snapshot checks visible restoration. It does not reset the internal
ID allocator. Shutdown uses \TestSource{shutdown/IsolatedShutdownTestSetup.java}
to verify the disposable target's process, port, and marker. External setup
provides a fresh target for each shutdown module.

\RunIn{Formats and invalid operations}
JSON/XML output checks cover resource and relationship reads, along with
resource creation and updates. Accept selects output separately from the
request Content-Type.
For example, \TestSource{todos/item/PutXmlResponseTest.java} compares the XML
update response with the stored todo. HEAD and OPTIONS check empty bodies.
Negative cases include unsupported verbs, invalid fields, and missing IDs.
Additionally,
\TestSource{categories/allCategories/PostMalformedExpectedTest.java} requires
400, unchanged data, and JSON Content-Type for the sampled malformed JSON/XML
bodies (BR07). Its companion \Path{PostMalformedObservedTest.java} checks the
actual \Path{text/html} label, parsed error messages, and unchanged state.
\TestSource{todos/item/DeleteTest.java} also deletes a todo twice and checks
that the second DELETE returns 404 without changing other data.

\RunIn{Execution and code structure}
We run these commands from \Path{testsuite/}, with Todo Manager started separately.
The first and third selections passed with Todo Manager 1.5.5 launched using
\Path{-version=4}. The second includes the failing expectations. The
\Path{expected-behavior} tag identifies the intended result, so tagged tests can
still pass. The default domain run includes these expectations and excludes shutdown.
\begin{lstlisting}
sh ./mvnw test '-Dgroups=!expected-behavior & !isolated-shutdown'
sh ./mvnw test -Dgroups=expected-behavior
sh ./mvnw test '-Dgroups=characterization | observed-behavior'
\end{lstlisting}
\Path{src/test/resources/junit-platform.properties} randomises both classes and
methods, with parallel execution disabled. Appending
\Path{'-Djunit.jupiter.execution.order.random.seed=1'} replays an order.
Fixture and parsing helpers avoid duplicated setup and parsing code.
\Path{categories/CategoryRelationshipDeletionTestSetup.java} separates scenario
creation, deletion, and verification into named helpers. Surefire writes XML
results to \Path{target/surefire-reports/} and prints class results to the console.
We also checked literal curl queries using \Path{--get --data-urlencode}.
% Exact commands, output, JSON/XML assertions and restoration:
% docs/project-part-a/section4-evidence/tag-verification/current-cli-query-evidence.json
Filtering by title and \Path{doneStatus} selected only the temporary todo.
An unmatched title returned an empty list. The before/after collections matched
after cleanup.
The seven Python scripts in \SessionScriptsFolder\ provide the exploratory
command-line demonstrations.
\par
\ifnum\value{page}>\numexpr\SuiteFirstPage+1\relax
  \PackageError{report}{Section 4 exceeds the requested two pages}%
    {Shorten Section 4 before recompiling.}
\fi

% ===================== BODY PAGE 5: SUBMISSION PACKAGE ================
\NewContentPage
\ReportSection{5. Submission Package and Reproduction}
\edef\PackageFirstPage{\number\value{page}}

\ReportSubsection{5.1 Package Structure}
Our source repository keeps the session records, demonstration scripts, and
Java suite together. The ZIP contents are shown below.
\begin{lstlisting}
README.md              Setup and reproduction instructions
Report & Documents/    Report, bug summary, capability inventory
Session Notes/          Eight session-note PDFs
Session Scripts/        Seven Python demonstrations
testsuite/              Java sources, Maven wrapper and configuration
Videos/                 Unit-test demonstration
\end{lstlisting}
Within \SuiteFolder, \Path{pom.xml} declares the build and dependencies,
\Path{.mvn/wrapper/} pins Maven, and \Path{src/test/} contains the tests,
shared support, and JUnit ordering settings. The README gives the full run
instructions, including preparation for the two isolated shutdown tests.

\ReportSubsection{5.2 Environment and Reproduction}
Use JDK 17 or newer. The wrapper downloads Maven 3.9.9 and the test dependencies.
Our recorded checks used Eclipse Adoptium JDK 21.0.12 on macOS 27.0.1 (aarch64).
Start the separately supplied Todo Manager 1.5.5 JAR before running Maven.
\Path{-version=4} selects the behaviour labelled \Path{v3} at startup.
\begin{lstlisting}
# Terminal 1: from the folder containing the supplied JAR
java -jar runTodoManagerRestAPI-1.5.5.jar -port=4567 -version=4
# Terminal 2: from the folder containing README.md
cd testsuite
sh ./mvnw clean test '-Djunit.jupiter.execution.order.random.seed=1'
sh ./mvnw clean test '-Djunit.jupiter.execution.order.random.seed=2'
\end{lstlisting}
In PowerShell, replace \Path{sh ./mvnw} with \Path{.\mvnw.cmd}.
Both classes and methods run in pseudorandom order. Reusing a seed and test
selection repeats that order. We keep expected-behaviour assertions enabled,
so API discrepancies produce failures. These domain runs exclude shutdown.
Section~4 gives the separate test selections.

Surefire prints results and writes XML to
\Path{testsuite/target/surefire-reports/}. Save each run's output before the next
\Path{clean}, which removes it. API documentation is available at
\url{http://localhost:4567/docs} while the service runs.

\RunIn{Session demonstrations}
The scripts use Python 3.9 or newer and the standard library. From the folder
containing the README, for example, run:
\begin{lstlisting}
python3 "Session Scripts/S6_category_relationships.py" --block 1
\end{lstlisting}
S5 and S6 accept \Path{--block N} to select one block. S7 and S8 also need the
JAR path as an argument and restart the service, resetting its data. Run scripts
and the suite separately so their changes do not overlap. S3's Postman requests
are in its session PDF, and the bug summary gives individual reproduction steps.
\par

% ===================== BODY PAGE 6: EXECUTION / VIDEO =================
\NewContentPage
% Check after flushing Section 5, including its final paragraph.
\ifnum\value{page}>\numexpr\PackageFirstPage+1\relax
  \PackageError{report}{Section 5 exceeds the requested one page}%
    {Shorten Section 5 before recompiling.}
\fi
\ReportSection{6. Findings of Unit-Test Execution}

\ReportSubsection{6.1 Recorded Results}
We ran the suite on October 20, 2026 using the unchanged repository revision
\Path{97208d8}. We used the environment in Section~5.2 and Todo Manager 1.5.5
with \Path{-version=4}, which starts the \Path{v3} profile.
For these checks, we started a separate process on port 4579.
Maven used \Path{-Dapi.baseUri=http://127.0.0.1:4579} and the seeds shown below.
The default runs kept our intended-behaviour assertions enabled.
% Fresh primary-source index and exact commands, including offline/cache flags:
% docs/project-part-a/section6-evidence/README.md
% full-seed1.json, full-seed2.json and their archived Surefire XML directories.
\begin{ReportTable}{@{}P{.29\linewidth}P{.12\linewidth}P{.12\linewidth}P{.12\linewidth}P{.10\linewidth}Y@{}}
  \toprule
  \textbf{Selection / seed} & \textbf{Executed} & \textbf{Passed} &
  \textbf{Failed} & \textbf{Errors} & \textbf{Skipped}\\
  \midrule
  Default / 1 & 917 & 753 & 164 & 0 & 0\\
  Default / 2 & 917 & 753 & 164 & 0 & 0\\
  Exclusion / 1 & 707 & 707 & 0 & 0 & 0\\
  Expected / 2 & 210 & 46 & 164 & 0 & 0\\
  Observed / 2 & 377 & 377 & 0 & 0 & 0\\
  \bottomrule
\end{ReportTable}
The table counts test invocations, including parameterized inputs.
The number after each selection is the pseudorandom ordering seed.
Default runs all domain tests from 256 classes. Exclusion removes tests tagged
\Path{expected-behavior}. Expected selects that tag. Observed selects
\Path{characterization} or \Path{observed-behavior}, which check the API
behaviour we recorded. The filtered rows rerun parts of Default, so their
counts are not additional coverage. Section~4.2 gives the commands.

The Failed column counts invocations with an unmet assertion.
We keep our intended expectations, so known discrepancies still count as failures.
Errors counts other problems running a test, such as uncaught exceptions.
Skipped counts disabled or aborted tests. All selected invocations produced
passes or assertion failures. Errors does not count HTTP error responses.
The POM excludes shutdown, so those two tests are not counted as skipped.
We saved console output and Surefire XML before later runs overwrote the reports.

Tests in Expected pass when the API meets our expectations. Of the 46 passing invocations,
44 checked rejection of unsupported methods and two checked task reads after
project deletion.

The same 164 invocations failed in both default runs. Their unmet expectations
include documented responses, relationship scoping and deletion, our policy of
hiding internal Java class names, and the unresolved BR30 supplied-ID hypothesis.
Therefore, 164 failed invocations do not mean 164 separate bugs or proven
documentation violations. Passing Observed checks reproduce the recorded
responses and stored effects. They do not show that a failed expectation has been fixed.

\RunIn{Malformed requests}
The category \Path{PostMalformedObservedTest} checks passed for the sampled
malformed JSON and XML bodies. They confirmed 400, parsed error messages, and
unchanged state. However, the matching media-type expectations failed because
the JSON-shaped error bodies were labelled \Path{text/html} (BR07).
Hiding the exposed Java class names is our separate quality expectation (BR08).

\RunIn{Isolated shutdown}
On separate disposable instances,
\TestSource{shutdown/GetObservedTest.java} passed its one invocation.
\Path{GetExpectedTest.java} failed its one invocation because the connection
closed without the documented HTTP 200. Both processes stopped and their ports
closed. Our first attempts also timed out during the process-exit check because
the external launcher had not collected the stopped process's exit status.
Correcting the launcher removed these timeouts. The test sources were unchanged.

\NewContentPage
\ReportSubsection{6.2 Readiness, Order Independence, and Side Effects}
\RunIn{Readiness}
We deliberately pointed \TestSource{todos/allTodos/GetTest.java} at the unused
port 4581 with \Path{-Dapi.baseUri=http://127.0.0.1:4581}. All 12 invocations
failed during setup with an instruction to start Todo Manager at that address.
There were no errors or skips. This checks the shared readiness failure through
one module.

\RunIn{Order and state}
Both seeds produced the same result for every invocation, although the console
showed different class orders. Seed~1 began with project task scope tests and
seed~2 with todo category scope tests. Classes and methods use pseudorandom
ordering. Parameterized inputs still follow their provider's order. This supports
independence only under the two tested orders.
Before/after snapshots matched for both full runs. Teardown checked restoration
after each test, including those that failed. The BR06 checks also compared
stored links and removed any created link before trying the next format. This
kept an earlier successful request from hiding a failed creation. These comparisons cover visible records and
relationship references. They do not restore the ID counter or verify hidden state.

\ReportSubsection{6.3 Bug Summary and Remaining Risks}
The following examples connect the execution results to the session findings
and the \BugSummaryRef. Their expected-behaviour checks failed, while the
companion observed-behaviour checks passed.
\begin{ReportTable}{@{}P{.13\linewidth}P{.44\linewidth}Y@{}}
  \toprule
  \textbf{Record} & \textbf{Execution finding and consequence} & \textbf{Session / module example}\\
  \midrule
  BR01, p.~1 & Missing parents returned unrelated targets with 200. The scope
  check does not require an undocumented 404 response. Clients could display
  another parent's records. &
  \NotesRef{6}, 9, p.~3. Category--project
  \Path{GetScopeExpectedTest} / \Path{GetScopeObservedTest}.\\
  BR04, p.~1 & Incoming-only links to a deleted category still returned its
  full record after direct GET returned 404. Clients could continue displaying
  deleted categories. &
  \NotesRef{6}, 19, p.~6. Category
  \Path{DeleteRelationshipsExpectedTest} / \Path{DeleteRelationshipsObservedTest}.\\
  BR06, p.~2 & Plain XML returned 404 for an existing target. JSON and CDATA
  linked that target with 201. Ordinary XML clients could fail to create the link. &
  \NotesRef{6}, 14, p.~4. Category--project
  \Path{PostXmlExpectedTest} / \Path{PostXmlObservedTest}.\\
  BR12, p.~3 & Existing category-link GET/HEAD/POST returned 404 instead of
  Swagger's 405. Clients could mistake an unsupported method for a missing link. &
  \NotesRef{6}, 18, p.~5. Individual category--project link
  \Path{GetExpectedTest} / \Path{GetObservedTest}.\\
  \bottomrule
\end{ReportTable}

\RunIn{Follow-up of the project reports}
The title-only PUT test in \TestSource{projects/item/PutObservedTest.java}
reset the description to empty and both status fields to false. Stored reads
showed that the task link was removed in both directions, while a second project's
link to the same todo remained. This answers the BR31--BR32 storage questions
for that request.
\TestSource{projects/item/tasks/PostExpectedTest.java} linked an existing todo
with 201 and verified both directions. It did not reproduce BR33 for its sampled
request. The difference from S3 still needs to be explained.

\RunIn{Overall assessment}
We would prioritise relationship scoping and deletion because clients could
receive unrelated or deleted records. However, the 33 bug records also contain
documentation gaps and policy questions. BR30 still needs an authoritative rule
on client-supplied IDs. The sampled project results do not establish every PUT
or link-creation case.

% ===================== BODY PAGE 7: AI DISCLOSURE =====================
\NewContentPage
\begingroup
% Give this prose section clear block spacing without changing other sections.
\setlength{\parskip}{0pt}
\newcommand{\DisclosureHeading}[1]{%
  \par\addvspace{12bp}%
  \noindent\textbf{#1:}\par\nobreak\vspace{4bp}}
\ReportSection{7. AI Assistance Disclosure}
\vspace{10bp}

\DisclosureHeading{AI Tools Used}
We used Claude Code and Codex mainly to help us with repetitive work in the unit
tests and project documents. We used these tools to set up the LaTeX format for our session notes,
bug summary, and report, and to turn our rough session notes into shorter, clearer
descriptions of what we tried and observed. We also used them to help draft and
revise parts of the report. We explored the API ourselves and recorded the core
observations in our session notes. For the unit tests, we used AI to help us map
those findings and our bug summary to the test suite. We chose what we
wanted to test and wrote pseudocode describing the steps and expected results.
AI then helped us convert this into Java unit tests and organise the suite around
the architecture we described, including separate tests for documented behaviour
and the actual behaviour we observed when the two differed.

We provided the API documentation, our existing session notes and scripts, and
the Java test sources to AI, which used these materials to generate the
capability inventory. We used the inventory to connect the
capabilities discussed in Sections~3.2 and~4 to their supporting demonstrations
and test modules.

\DisclosureHeading{Verification of AI Output}
To check that the AI-assisted code tested what we intended, we read through the
unit tests and compared their setup, requests, and assertions with our pseudocode,
the observations we recorded during our exploration, and the REST API documentation.
We also ran the tests and checked
the results against the behaviour we expected. For cases where the API behaved
differently from the documentation, we checked both the assertions for the
documented behaviour and the separate tests for the behaviour recorded in our
session notes. For the written documents, we reviewed the revised wording against
our original notes to make sure that shortening the text did not change what we
had tried or observed. We also checked the tables, images, and LaTeX formatting
to make sure that the documents displayed properly and compiled without errors.

\DisclosureHeading{Strengths and Limitations}
AI helped us save time by applying the same structure to similar test cases and
formatting our documents consistently in LaTeX. Once we had worked out the
observations, expected results, and test steps, it helped us repeat that
implementation and formatting work across the project. However, it could
also misunderstand an observation or generate a test with an incorrect
expectation. A test compiling or passing does not show that it checks the right
thing, so we still needed to read the code and check what each assertion was
actually confirming. The same applied to our session notes: a clearer description
still needed to match the original observation.

\DisclosureHeading{Conclusion}
Overall, we used AI mainly to handle the repetitive work involved in implementing
our test ideas and preparing the project documents. We explored the API, made the
core observations, and decided what needed to be tested. AI helped us turn those
findings and our pseudocode into Java tests and readable documents, and we were
still responsible for reviewing the final code, checking the results, and making
sure that the written descriptions matched our work.
\par
\endgroup

% ===================== BODY PAGE 8: REFERENCES ========================
\NewContentPage
\ReportSection{8. References and Supporting Artifacts}
The artifact links use paths within the extracted submission ZIP, with this
report in \Path{Report & Documents/}. Each session reference below identifies
one PDF in \Path{Session Notes/}.
% List labels supply brackets; cross-references return only the number.
\begin{enumerate}[label={[\arabic*]},ref=\arabic*,leftmargin=2em,itemsep=3bp,topsep=3bp]
  \item\label{ref:api} Todo Manager 1.5.5 REST API documentation, served locally
    while the application is running:
    \url{http://localhost:4567/docs} and
    \url{http://localhost:4567/docs/swagger}.
  \item\label{ref:session-1} Jeffrey Lim. S1: Todo CRUD.\\
    \ArtifactFile{../Session Notes/SessionNotes_S1_Jeffrey_Lim.pdf}
      {SessionNotes_S1_Jeffrey_Lim.pdf}.
  \item\label{ref:session-2} Jeffrey Lim. S2: Todo relationships.\\
    \ArtifactFile{../Session Notes/SessionNotes_S2_Jeffrey_Lim.pdf}
      {SessionNotes_S2_Jeffrey_Lim.pdf}.
  \item\label{ref:session-3} Mahan Maalekian. S3: Project CRUD.\\
    \ArtifactFile{../Session Notes/SesssionNotes_S3_Mahan_Maalekian.pdf}
      {SesssionNotes_S3_Mahan_Maalekian.pdf}.
  \item\label{ref:session-4} James Helou. S4: Project relationships.\\
    \ArtifactFile{../Session Notes/SessionNotes_S4_James_Helou_261177052.pdf}
      {SessionNotes_S4_James_Helou_261177052.pdf}.
  \item\label{ref:session-5} Matthew Petruzziello. S5: Category CRUD.\\
    \ArtifactFile{../Session Notes/SessionNotes_S5_Matthew_Petruzziello.pdf}
      {SessionNotes_S5_Matthew_Petruzziello.pdf}.
  \item\label{ref:session-6} Matthew Petruzziello. S6: Category relationships.\\
    \ArtifactFile{../Session Notes/SessionNotes_S6_Matthew_Petruzziello.pdf}
      {SessionNotes_S6_Matthew_Petruzziello.pdf}.
  \item\label{ref:session-7} James Helou. S7: Interoperability and deletion.\\
    \ArtifactFile{../Session Notes/SessionNotes_S7_James_Helou_261177052.pdf}
      {SessionNotes_S7_James_Helou_261177052.pdf}.
  \item\label{ref:session-8} James Helou. S8: Malformed and invalid inputs.\\
    \ArtifactFile{../Session Notes/SessionNotes_S8_James_Helou_261177052.pdf}
      {SessionNotes_S8_James_Helou_261177052.pdf}.
  \item\label{ref:bug-summary} Group 3. Bug summary, BR01--BR33.\\
    \BugSummaryLink, in \Path{Report & Documents/}.
  \item\label{ref:video} Group 3. Unit-test demonstration.\\
    \UnitTestVideoLink.
  \item\label{ref:capabilities} Group 3. Complete capability and demonstration inventory.\\
    \CapabilityInventoryLink, in \Path{Report & Documents/}.
    Lists documented methods, discovered behaviours, session/script evidence,
    and the Java modules used in Section~4.
\end{enumerate}

% FINAL DRAFT AUDIT (not printed):
% Resolve every Draft/Prompt; remove DraftNotice; replace figure placeholders.
% Check claims against actual notes, scripts, tests, bug records, and video.
% Keep the final report 5--10 pages. Verify actual submission timing separately.
% Move only the final PDF into ecse-429-projectA when the user is ready.
\end{document}
